% =================================================
% 性学系列 · 男性性健康
% 共享基座见 common.tex
% =================================================
\documentclass[12pt,UTF8,openany]{ctexbook}
% =================================================
% common.tex —— 性学系列 共享基座
% 导言区 / 配色 / TikZ 库 / ctexset / C 列型 / 字体 / 书眉页码
% 各册 % =================================================
% common.tex —— 性学系列 共享基座
% 导言区 / 配色 / TikZ 库 / ctexset / C 列型 / 字体 / 书眉页码
% 各册 % =================================================
% common.tex —— 性学系列 共享基座
% 导言区 / 配色 / TikZ 库 / ctexset / C 列型 / 字体 / 书眉页码
% 各册 \input{common.tex} 引用
% =================================================
\usepackage[top=2.5cm,bottom=2.5cm,left=2.6cm,right=2.6cm,headheight=2.0cm]{geometry}
\usepackage{amsmath}
\usepackage{amssymb}
\usepackage{array}
\usepackage{tabularx}
\usepackage{multirow}
\usepackage{booktabs}
\usepackage{enumitem}
\usepackage{xcolor}
\usepackage{graphicx}
\usepackage{tikz}
\usetikzlibrary{positioning,calc,arrows.meta,shapes.geometric,matrix}
\definecolor{myblue}{HTML}{185FA5}
\definecolor{myblueL}{HTML}{E6F1FB}
\definecolor{mygreen}{HTML}{3B6D11}
\definecolor{mygreenL}{HTML}{C0DD97}
\definecolor{myteal}{HTML}{1D9E75}
\definecolor{myamber}{HTML}{854F0B}
\definecolor{myamberL}{HTML}{FAEEDA}
\definecolor{myred}{HTML}{A32D2D}
\definecolor{myredL}{HTML}{FCEBEB}
\definecolor{mypurple}{HTML}{534AB7}
\definecolor{mypurpleL}{HTML}{EEEDFE}
\definecolor{mycoral}{HTML}{D85A30}
\definecolor{mycoralL}{HTML}{FAECE7}
\definecolor{mypink}{HTML}{D4537E}
\definecolor{mypinkL}{HTML}{F4C0D1}
\definecolor{mypinkM}{HTML}{ED93B1}
\definecolor{mygray}{HTML}{888780}

% 字体
\setCJKfamilyfont{hei}{SimHei}
\newcommand{\myhei}{\CJKfamily{hei}}

% 书眉与页码（印刷友好：中文页眉 + 居中页码）
\usepackage{fancyhdr}
\pagestyle{fancy}
\fancyhf{}
\fancyhead[L]{\small\myhei 性学系列}
\fancyhead[R]{\small\leftmark}
\fancyfoot[C]{\small\thepage}
\renewcommand{\headrulewidth}{0.4pt}
\fancypagestyle{plain}{\fancyhf{}\fancyfoot[C]{\small\thepage}\renewcommand{\headrulewidth}{0pt}}
\fancypagestyle{empty}{\fancyhf{}\renewcommand{\headrulewidth}{0pt}}

% 行距（零基础易读，略松）
\linespread{1.3}

% 超链接：印刷友好（黑色，电子版仍可点击）
\usepackage[colorlinks,linkcolor=black,citecolor=black,urlcolor=black]{hyperref}

\ctexset{
  chapter/name={第,章},
  chapter/number={\arabic{chapter}},
  section/format=\Large\bfseries,
}

\newcolumntype{C}{>{\centering\arraybackslash}X}
 引用
% =================================================
\usepackage[top=2.5cm,bottom=2.5cm,left=2.6cm,right=2.6cm,headheight=2.0cm]{geometry}
\usepackage{amsmath}
\usepackage{amssymb}
\usepackage{array}
\usepackage{tabularx}
\usepackage{multirow}
\usepackage{booktabs}
\usepackage{enumitem}
\usepackage{xcolor}
\usepackage{graphicx}
\usepackage{tikz}
\usetikzlibrary{positioning,calc,arrows.meta,shapes.geometric,matrix}
\definecolor{myblue}{HTML}{185FA5}
\definecolor{myblueL}{HTML}{E6F1FB}
\definecolor{mygreen}{HTML}{3B6D11}
\definecolor{mygreenL}{HTML}{C0DD97}
\definecolor{myteal}{HTML}{1D9E75}
\definecolor{myamber}{HTML}{854F0B}
\definecolor{myamberL}{HTML}{FAEEDA}
\definecolor{myred}{HTML}{A32D2D}
\definecolor{myredL}{HTML}{FCEBEB}
\definecolor{mypurple}{HTML}{534AB7}
\definecolor{mypurpleL}{HTML}{EEEDFE}
\definecolor{mycoral}{HTML}{D85A30}
\definecolor{mycoralL}{HTML}{FAECE7}
\definecolor{mypink}{HTML}{D4537E}
\definecolor{mypinkL}{HTML}{F4C0D1}
\definecolor{mypinkM}{HTML}{ED93B1}
\definecolor{mygray}{HTML}{888780}

% 字体
\setCJKfamilyfont{hei}{SimHei}
\newcommand{\myhei}{\CJKfamily{hei}}

% 书眉与页码（印刷友好：中文页眉 + 居中页码）
\usepackage{fancyhdr}
\pagestyle{fancy}
\fancyhf{}
\fancyhead[L]{\small\myhei 性学系列}
\fancyhead[R]{\small\leftmark}
\fancyfoot[C]{\small\thepage}
\renewcommand{\headrulewidth}{0.4pt}
\fancypagestyle{plain}{\fancyhf{}\fancyfoot[C]{\small\thepage}\renewcommand{\headrulewidth}{0pt}}
\fancypagestyle{empty}{\fancyhf{}\renewcommand{\headrulewidth}{0pt}}

% 行距（零基础易读，略松）
\linespread{1.3}

% 超链接：印刷友好（黑色，电子版仍可点击）
\usepackage[colorlinks,linkcolor=black,citecolor=black,urlcolor=black]{hyperref}

\ctexset{
  chapter/name={第,章},
  chapter/number={\arabic{chapter}},
  section/format=\Large\bfseries,
}

\newcolumntype{C}{>{\centering\arraybackslash}X}
 引用
% =================================================
\usepackage[top=2.5cm,bottom=2.5cm,left=2.6cm,right=2.6cm,headheight=2.0cm]{geometry}
\usepackage{amsmath}
\usepackage{amssymb}
\usepackage{array}
\usepackage{tabularx}
\usepackage{multirow}
\usepackage{booktabs}
\usepackage{enumitem}
\usepackage{xcolor}
\usepackage{graphicx}
\usepackage{tikz}
\usetikzlibrary{positioning,calc,arrows.meta,shapes.geometric,matrix}
\definecolor{myblue}{HTML}{185FA5}
\definecolor{myblueL}{HTML}{E6F1FB}
\definecolor{mygreen}{HTML}{3B6D11}
\definecolor{mygreenL}{HTML}{C0DD97}
\definecolor{myteal}{HTML}{1D9E75}
\definecolor{myamber}{HTML}{854F0B}
\definecolor{myamberL}{HTML}{FAEEDA}
\definecolor{myred}{HTML}{A32D2D}
\definecolor{myredL}{HTML}{FCEBEB}
\definecolor{mypurple}{HTML}{534AB7}
\definecolor{mypurpleL}{HTML}{EEEDFE}
\definecolor{mycoral}{HTML}{D85A30}
\definecolor{mycoralL}{HTML}{FAECE7}
\definecolor{mypink}{HTML}{D4537E}
\definecolor{mypinkL}{HTML}{F4C0D1}
\definecolor{mypinkM}{HTML}{ED93B1}
\definecolor{mygray}{HTML}{888780}

% 字体
\setCJKfamilyfont{hei}{SimHei}
\newcommand{\myhei}{\CJKfamily{hei}}

% 书眉与页码（印刷友好：中文页眉 + 居中页码）
\usepackage{fancyhdr}
\pagestyle{fancy}
\fancyhf{}
\fancyhead[L]{\small\myhei 性学系列}
\fancyhead[R]{\small\leftmark}
\fancyfoot[C]{\small\thepage}
\renewcommand{\headrulewidth}{0.4pt}
\fancypagestyle{plain}{\fancyhf{}\fancyfoot[C]{\small\thepage}\renewcommand{\headrulewidth}{0pt}}
\fancypagestyle{empty}{\fancyhf{}\renewcommand{\headrulewidth}{0pt}}

% 行距（零基础易读，略松）
\linespread{1.3}

% 超链接：印刷友好（黑色，电子版仍可点击）
\usepackage[colorlinks,linkcolor=black,citecolor=black,urlcolor=black]{hyperref}

\ctexset{
  chapter/name={第,章},
  chapter/number={\arabic{chapter}},
  section/format=\Large\bfseries,
}

\newcolumntype{C}{>{\centering\arraybackslash}X}


\begin{document}

% ===== 封面 =====
\begin{titlepage}
\centering
\vspace*{1.5cm}
\begin{tikzpicture}
\draw[myblue,line width=2pt] (0,0) circle (1.2);
\node[text=white] at (0,0) {\myhei\zihao{3}性};
\end{tikzpicture}
\vspace{0.6cm}
\noindent\rule{\textwidth}{0.4pt}
\vspace{0.5cm}
\myhei\zihao{5} 性学系列\,·\,家庭健康科普\,·\,全五册
\vspace{2.5cm}
\myhei\zihao{0} 男性性健康\par
\vspace{1.2cm}
\itshape\zihao{3} 男性身体的专项知识\par
\vspace{5.2cm}
\noindent\rule{\textwidth}{0.4pt}
\vspace{0.5cm}
\zihao{5} 性学系列\,·\,\today
\end{titlepage}

% ===== 免责声明 =====
\begin{center}\zihao{5}\itshape
本书为家庭健康科普读物，内容力求科学严谨，但\textbf{不替代专业医生的诊断与治疗}。\\
涉及个人健康问题，请以正规医疗机构的检查与建议为准。\\
全书以安全、自愿、相互尊重为底色，尊重多元与个体差异。\par\end{center}
\clearpage

% ===== 如何使用本书 =====
\chapter*{如何使用本书}
\addcontentsline{toc}{chapter}{如何使用本书}
\noindent\textbf{性学系列共五册：}
\begin{itemize}[leftmargin=2em]
\item \textbf{A 性健康基础} —— 认知、同意、防护、STI、避孕总论（建议先读）
\item \textbf{B 亲密实践与技巧} —— 前戏、体位、高潮引导、呼吸与盆底
\item \textbf{C 男性性健康} —— 男性解剖、ED/自慰、绝育、不育
\item \textbf{D 女性性健康} —— 女性解剖、妇科、不孕、人流、避孕、孕期/更年期
\item \textbf{E BDSM 与 Kink} —— 安全框架下的亲密表达
\end{itemize}
\vspace{0.6em}
\noindent\textbf{建议阅读路径（零基础）：} 先通读 A 建立认知与安全底线 → 按需读 C 或 D 了解自己的身体 → 再读 B 学技巧 → 兴趣到了再读 E。每章末附有本章小结，便于快速回顾。
\clearpage

% ===== 目录 =====
\tableofcontents
\clearpage
\chapter{男性视角的补充：同样重要的一半}
\label{chap:male}

前面各章更多是"如何保护自己"的口吻。但性是两个人的事，
男性有对等的权利、责任与困扰。这一章从男性角度补齐。

\section{同意是双向的：你也有权说不}

\begin{itemize}[leftmargin=2em]
  \item 前面讲"学会说不"，对男性同样成立。\textbf{对方有同意权，你也有}。
  \item 不必因为"是男人"就必须主动、必须想、必须持续。没准备好就说"我今天不想"，不必解释。
  \item 被对方施压（"你不配合就是不爱我"）和你施压一样，都是越界。
  \item \textbf{别把"不抗拒"当成"同意"}——明确、清醒、可撤回的肯定，才是同意。
\end{itemize}

\section{男性也会成为受害者}

这不是小事。WHO 数据显示，约\textbf{1/5 的女性、1/7 的男性}报告童年遭受过性侵。
男性遭遇性骚扰、性侵、亲密关系暴力的比例不低，但因为"男人就该扛着"的观念，
很多人选择沉默。

\begin{itemize}[leftmargin=2em]
  \item 遭遇侵害不是你的错，与"强不强壮"无关。
  \item 求助渠道对男性同样开放：拨打 110、12320，或联系医院、可信任的人。
  \item WHO 明确指出：僵化的"男性气概"规范会损害男性健康——鼓励冒险、酗酒、不愿求助。
        主动求助是勇敢，不是软弱。
\end{itemize}

\section{常见的身体困扰：多见、可治、勿盲目进补}

\subsection*{勃起功能障碍（ED）：全身健康的信号}

\begin{itemize}[leftmargin=2em]
  \item 不同研究数字差异很大（如波士顿 MMAS 研究 40--70 岁约 52\%，
        科隆研究 30--80 岁约 19.2\%），共同结论是\textbf{随年龄上升}。
  \item 关键：ED 常与糖尿病、高血压、血脂异常等代谢/心血管问题相关。
        EAU 指南要求在诊治时评估血糖、血脂和睾酮——\textbf{它可能是身体给出的早期警报}，
        值得正规就医，而不是买补药。
  \item 治疗是一套成熟方案：先改善生活方式（运动、减重、控压控糖），
        一线药物为 PDE5 抑制剂，均需在医生指导下使用。
\end{itemize}

\textbf{一线药两种主力——按需选哪种（均需处方）：} 最常用的两类 PDE5 抑制剂，关键差别在\textbf{半衰期长短}，决定了"按需用"还是"每天吃"。

\begin{tabularx}{\textwidth}{l X X}
\toprule
& 西地那非（品牌：万艾可/伟哥） & 他达拉非（品牌：希爱力） \\
\midrule
半衰期 & 约 4--5 小时 & 约 17.5 小时（长 3 倍以上） \\
有效窗口 & 服后约 0.5--4 小时 & 约 24 小时（"周末药/假日药"，一天内随时可用） \\
每日低剂量 & 不适用 & \textbf{支持}（5--10\,mg 每日，适合规律性生活、心理性 ED） \\
受进食影响 & 高脂餐使起效延迟 & 基本不受进食影响 \\
常见副作用 & 头痛、面部潮红、视物偏蓝 & 头痛、肌痛/背痛（持续较久）、鼻塞 \\
\midrule
共同禁忌 & \multicolumn{2}{c}{\textbf{正在服硝酸酯类（如硝酸甘油）者禁用}；严重心血管病慎用} \\
\bottomrule
\end{tabularx}

\textbf{如何选择（需遵医嘱）：} 性生活不规律、想"即用即起"——偏西地那非这类短效药；生活较规律、不想每次提前算时间、或心理因素为主——偏他达拉非每日低剂量。两者都\textbf{需性刺激才起效}（并非"吃了就自动勃起"）、都\textbf{不预防性病}；按需服用者不可连续无间断服用。

\textbf{硝酸酯类 + PDE5 抑制剂——唯一"可预防但常被忽视"的致命组合：} 两者都扩张血管、降低血压，\textbf{叠加后血压骤降}（可低至 70/40\,mmHg 以下），表现为头晕、视物模糊、心跳快、出冷汗，严重时\textbf{晕厥、心梗、猝死}。这不是"不舒服"，是急性血流动力学崩溃。

\begin{tabularx}{\textwidth}{l X}
\toprule
\textbf{关键} & \textbf{说明} \\
\midrule
\textbf{哪些药算"硝酸酯类"} & 共 4 大类：\textbf{硝酸甘油}、\textbf{硝酸异山梨酯}、\textbf{单硝酸异山梨酯}、\textbf{硝普钠}。具体商品名识别见下表——很多患者只把"硝酸甘油"当硝酸酯，拿着"依姆多/异乐定/消心痛"就以为无关，\textbf{看通用名带"硝酸"或"异山梨酯"、或商品名在下表内就警惕}。 \\
\textbf{安全间隔} & 半衰期决定。西地那非（4--5\,h）与硝酸酯间隔\textbf{ $\ge$24\,h}；他达拉非（17.5\,h）间隔\textbf{ $\ge$48\,h}。\textbf{反过来也成立}：用了他达拉非 48\,h 内，\textbf{心绞痛发作不能直接舌下含硝酸甘油}——这是常被忽略的"反向禁忌"。 \\
\textbf{正在用硝酸酯类怎么办} & \textbf{不能为"能勃起"就停硝酸酯}（硝酸酯是保命药）。找 ED 的替代方案：\textbf{真空负压装置}（无药物相互作用）、\textbf{海绵体注射}（前列地尔/卡前列素，局部作用）、\textbf{尿道栓剂}（muse）；严重或长期可考虑\textbf{阴茎假体植入}。亚甲蓝（methylene blue）可替代部分场景，但它本身也降压、且是强 CYP 抑制剂，\textbf{须医生评估}。 \\
\textbf{紧急处理（已误服、出现症状）} & \textbf{立即平躺、抬高双腿、不要站起}（防体位性低血压加重）；如血压 $<$90/60\,mmHg 或意识模糊，\textbf{立即急诊}。不要"再等一会看看"——低血压性晕厥可致跌倒外伤、脑供血不足。 \\
\textbf{就医沟通（最关键）} & 看\textbf{任何}医生（心内科、心外科、急诊、麻醉）都要\textbf{主动说"我在吃他达拉非/西地那非"}——别让医生在不知情的情况下给你开硝酸甘油。反过来，看男科/泌尿科时也要主动说\textbf{"我在用硝酸酯类"}。 \\
\bottomrule
\end{tabularx}

\textbf{硝酸酯类药物识别：}

\begin{tabularx}{\textwidth}{p{3.6cm} X}
\toprule
\textbf{通用名（学名）} & \textbf{常见商品名 / 别名 · 剂型 · 用途} \\
\midrule
硝酸甘油（三硝酸甘油酯） & \textbf{保欣宁、乃才朗、礼顿、疗痛脉、永保心灵}。剂型：舌下片、气雾剂（喷雾）、贴片、注射液。用途：急性心绞痛发作（救心），起效最快、数分钟。 \\
硝酸异山梨酯（二硝酸异山梨酯） & \textbf{消心痛、异舒吉、长效易顺脉、ISOKET、SORBITRATE}。剂型：口含片、片剂、缓释片、注射液。用途：心绞痛治疗与预防。 \\
单硝酸异山梨酯 & \textbf{依姆多(Imdur)、异乐定、欣康、爱倍、诺可达、鲁南欣康、安心脉}。剂型：片剂、缓释片、缓释胶囊、注射液。用途：心绞痛长期治疗/预防（长效）。 \\
硝普钠 & 多为通用名（少见商品名）。剂型：注射液（住院静脉滴注）。用途：高血压急症、急性心衰——\textbf{虽名不带"硝酸酯"，但同属一氧化氮供体，与 PDE5 抑制剂同样致命}。 \\
\bottomrule
\end{tabularx}

\textbf{快速识别：} 通用名带"硝酸"或"异山梨酯"，或商品名是"消心痛/依姆多/异乐定/欣康/诺可达/爱倍/鲁南欣康/安心脉/乃才朗/保欣宁"之一 $\to$ 一律当硝酸酯对待，\textbf{吃他达拉非/西地那非前先问医生}。

\textbf{核心原则：} 硝酸酯是"保命药"，PDE5 是"提升生活质量的药"——两者不能共存，\textbf{保命药永远优先}。正在用硝酸酯类的男性，ED 的解决方案不是"停硝酸酯改吃他达拉非"，而是\textbf{换一种不依赖全身血管的 ED 治疗}。

\subsection*{早泄：先判断是否属于疾病}

\begin{itemize}[leftmargin=2em]
  \item ICD-11 已把它改称"早期射精"：指插入前或插入后极短时间内射精、
        几乎无控制感、持续数月并带来明显困扰。
  \item EAU 强调\textbf{不能只看时间}，还要看控制感、是否普遍、是否造成困扰。
        "正常变异"（时快时慢）和"主观性早泄"（时间其实正常但自认为快）都不算真正的疾病。
  \item 分\textbf{原发}（一直如此）与\textbf{继发}（后来才出现，可能与前列腺、激素、药物或心理因素有关）——继发的更要先找原因。
  \item \textbf{心态}：别把它当"个人缺陷"——它是\textbf{常见、可管理}的，伴侣的理解与沟通本身就是治疗的一部分。
  \item 真正需要时，行为训练、盆底肌锻炼、药物等都是有效手段，去泌尿外科/男科评估即可，\textbf{别独自硬扛或乱买补药}。
\end{itemize}

\textbf{治疗路径（从简到繁，多数轻度不需要药）：}

\begin{itemize}[leftmargin=2em,itemsep=0.3em]
  \item \textbf{行为法·停-起始法}：刺激到将射未射（约七八成兴奋）时停住，待射精冲动消退后再继续，如此重复数轮、最后一轮再完成。目的是让大脑重新识别"临界点"、建立"刹车"反射。自慰时即可练习，不必等到同房。
  \item \textbf{行为法·挤压法}：接近高潮时，用指腹轻压龟头下方（系带处）数秒，直到射精冲动消退。
  \item \textbf{盆底肌（PC 肌）锻炼}：\textbf{先找到它}——排尿中途能"刹车"的那组肌肉就是（此法只用来\textbf{定位}，找到后就改用日常状态下的收缩训练；\textbf{长期靠中断排尿来练}可能引起排尿不畅或残余尿）。收缩保持约 5 秒、放松约 5 秒为一组，每组 10 次、每天 2--3 组，从慢开始、逐步加快，坚持数周见效；注意\textbf{单独收缩该肌、不憋气、不夹臀}。它同时改善勃起硬度与射精控制。
  \item \textbf{局部麻醉剂（延时喷剂/药膏）}：含利多卡因/丙胺卡因等，降低龟头敏感度，无系统性副作用；可能波及伴侣敏感度，注意用量、吸收后再行。
  \item \textbf{处方药（需医生评估）}：\textbf{达泊西汀}是目前唯一获批按需用于早泄的口服药，性生活前 1--3 小时服用、按需而非每日；另有部分 SSRI（如帕罗西汀/舍曲林）可超说明书用于早泄，需\textbf{每日服用、起效较慢（1--2 周}），有性欲下降、嗜睡等副作用。\textbf{切勿自行抓药，或与正在服用的抗抑郁药混用。}
\end{itemize}

\textbf{达泊西汀的副作用与应对（多为轻度，1--2 次后常减轻）：} 常见头晕、恶心、嗜睡、头痛、腹泻、视物模糊。应对：随餐或进食后服用可减轻恶心；起身/站起时放慢动作，防体位性头晕；\textbf{服药当日避免饮酒}（会加重头晕/低血压）；\textbf{正在服硝酸酯类（如硝酸甘油）或强 CYP3D4 抑制剂者禁用}，严重肝肾功能不全者慎用——这些禁忌需医生把关，\textbf{别自行抓药或加量}。

\textbf{与他达拉非（治 ED）副作用对比：} 两者总体都是\textbf{轻--中度、多自限}，量级相当（约两成到三成的人出现至少一种主要症状）。差别在"谱系"与"持续时间"：\textbf{达泊西汀}偏向中枢与胃肠（\textbf{恶心最突出}，其次头晕、嗜睡），因半衰期短而\textbf{来得快去得快}（数小时内消退），最需警惕\textbf{体位性低血压/晕厥}及与抗抑郁药合用诱发 5-羟色胺综合征；\textbf{他达拉非}偏向血管与肌肉（\textbf{头痛}最常见、\textbf{背痛/肌痛}为其特色），因半衰期长，\textbf{副作用（尤其肌痛）可持续 24\,h 以上}，最需警惕与\textbf{硝酸酯类}合用引发的严重低血压。两药\textbf{合并使用}（PE+ED）时，达泊西汀的头晕与他达拉非的低血压效应会\textbf{叠加}——更须避免同日饮酒、起身放慢、不与硝酸酯/强 CYP3D4 抑制剂同服，首次建议错开并观察自身反应。

\textbf{先排查可逆病因}：慢性前列腺炎、长期大量饮酒、焦虑与"表现压力"都可表现为早泄——\textbf{治原发病}往往比单纯"延时"更根本（本册“男性不育”章论及前列腺、“表现焦虑”一节论及压力）。

\textbf{如果同时有勃起困难（ED）：} 临床上"既快又硬不起来"并不少见。处理上\textbf{通常先评估并处理 ED}：一是先查器质性原因（糖尿病/心血管/血脂——ED 常是全身健康信号）；二是勃起信心恢复后，因紧张导致的早泄常随之改善；三是确需药物时，PDE5 抑制剂（治 ED）与达泊西汀（治早泄）\textbf{可联合使用}，但\textbf{都需医生评估}（服心血管药者、硝酸酯类禁用 PDE5i 尤须注意）。详见本册“勃起功能障碍”一节。

\subsection*{对“壮阳保健品”保持警惕}

来路不明的补药、偏方成分不明，可能有害且耽误真正病因的诊治。有困扰就去看正规医院。

\section{男性在避孕与防病中的责任}

\begin{itemize}[leftmargin=2em]
  \item \textbf{安全套（避孕套）}：男性主导使用的工具，主动准备好、正确佩戴，是责任也是保护自己。
  \item \textbf{输精管结扎}：高效、永久性的男性避孕方式。CDC 数据：典型使用第一年失败率约
        0.15\%；术后 8--16 周需做精液检查确认无精子，确认前仍需用屏障避孕。
        它\textbf{不预防性病}。后悔率约 5\%，决定前要充分咨询。
  \item \textbf{HPV 疫苗男性也要打}：CDC 建议男女都在 11--12 岁常规接种、可 9 岁起，
        未接种者 26 岁前补种，27--45 岁与医生共同决策。可预防生殖器疣、肛门癌等。
  \item \textbf{乙肝疫苗}：有多个性伴或接受性病诊疗者应接种。
  \item 定期做性病筛查，同样是对自己和伴侣负责。
\end{itemize}

\section{“表现焦虑”：多为习得而非天生}

\begin{itemize}[leftmargin=2em]
  \item 色情内容里的时长、强度、反应都是表演，拿它当标准必然焦虑。
  \item 一次性生活"不理想"非常常见——疲劳、压力、饮酒、情绪都会影响。
        偶尔失常不等于"不行了"。
  \item 越关注"表现"，越容易真的出问题（焦虑本身就是 ED 的常见原因）。
        把注意力放回连接、沟通和感受，反而更好。
  \item 真正需要警惕的是问题\textbf{持续出现}并带来困扰——那就去就医，而不是独自硬扛。
\end{itemize}

\section{本章小结}

\begin{enumerate}[leftmargin=2em,itemsep=0.3em]
  \item 我尊重对方的同意，也保有自己说"不"的权利。
  \item 我知道男性也可能成为受害者，遇到侵害会求助而不是硬扛。
  \item 我明白 ED/早泄常见且可治疗，ED 还可能是全身健康信号，会去正规医院而非买补药。
  \item 我会主动准备安全套，知道输精管结扎与 HPV 疫苗等选项。
  \item 我不拿色情内容当标准，不把一次性失常当成"不行"。
  \item 我不迷信"男人气概"，会表达情绪、会求助。
\end{enumerate}

\chapter{当伴侣遇到勃起困难：支持、沟通与自慰的正常认知}
\label{chap:partner-ed}

ED（勃起功能障碍）前面从"当事人自己怎么办"讲过，本章换一个很少被正面讲的视角——\textbf{伴侣如何面对和支持}，并顺带澄清一个常被病理化的话题：\textbf{自慰}。核心信息有三：ED 常常是\textbf{健康信号}而非"面子问题"；现代\textbf{有一线药}可用；自慰是\textbf{正常的性表达}。

\section{术语界定：ED 而非“阳痿”}
叫\textbf{勃起功能障碍（ED）}，不叫"阳痿"。"阳痿"带强烈污名和羞辱，本身就会加重焦虑。大多数男性一生中\textbf{偶尔都会遇到}，\textbf{持续 3 个月以上}才需系统评估。名字一改，求助门槛就低了。

\section{成因与心血管预警（ED 常常是健康信号）}
血管问题最常见；糖尿病、高血压、吸烟、高血脂、部分降压/抗抑郁药都可能参与；心理因素（紧张、表现焦虑、关系紧张）常与生理因素\textbf{叠加}，不是"二选一"。
\textbf{关键点：ED 常常是全身血管健康的早期信号}——它可能与心血管疾病独立相关，甚至可成为动脉粥样硬化的\textbf{首个临床表现}。出现 ED 该评估血压、血糖、血脂、吸烟——这不是"面子问题"，是\textbf{健康预警}。

\section{常见误区辨析}
\begin{tabularx}{\textwidth}{l X}
\toprule
\textbf{常见误区} & \textbf{科学认识} \\
\midrule
"阳痿"是固定叫法 & 现代用"勃起功能障碍（ED）"，去污名化有助求助。 \\
"主要靠心理/肾亏调理" & 多为血管/混合因素；\textbf{PDE5 抑制剂（西地那非/他达拉非）是一线}，强推荐。 \\
"ED 就是个人问题" & 常是\textbf{心血管健康的预警信号}，需评估整体风险。 \\
"自慰 = 婚姻不满足/精神不忠" & 自慰是\textbf{正常、普遍、无害}的性表达，不等于不忠；只有强迫性、明显影响生活才需关注。 \\
"只能忍" & 有药、有方法；\textbf{早评估、早干预}，不必"忍多年"。 \\
\bottomrule
\end{tabularx}

\section{伴侣的支持要点}
\begin{itemize}[leftmargin=2em,itemsep=0.3em]
  \item \textbf{要}：不指责、一起就医、把 ED 当\textbf{健康问题}（不是"爱不爱"的证明）、接受\textbf{亲密不等于插入}（拥抱/爱抚/相互自慰都算）。
  \item \textbf{不要}：不当场追问施压、不把它解读成"嫌弃我"、不替对方隐瞒回避。
  \item \textbf{心理冲击是双向的}：男方可能担心"我不是个失败的男人吗？她会离开我吗？"；女方可能误以为"是我没吸引力了/他有外遇了"——\textbf{隐瞒与回避}反而让关系更冷，\textbf{一起面对}才关键。
\end{itemize}

\section{现代治疗：PDE5 抑制剂为一线}
\begin{itemize}[leftmargin=2em,itemsep=0.3em]
  \item \textbf{一线}：PDE5 抑制剂（西地那非/他达拉非，\textbf{需处方}；\textbf{正在服硝酸酯类心血管药者禁用}）。
  \item \textbf{生活方式}：戒烟、运动、控糖控压控脂——对血管性 ED 有实质帮助。
  \item \textbf{二线}：无效再考虑真空装置、海绵体注射、假体手术。
  \item \textbf{心理/关系}：性治疗、伴侣咨询；ED 与焦虑常互为因果。
  \item \textbf{别漏查}：激素（睾酮/甲状腺/催乳素）、心血管风险。
\end{itemize}

\section{自慰的正常化}
自慰是\textbf{正常、普遍、无害}的性表达，不需要羞耻，也\textbf{不等于对伴侣不忠}。它是\textbf{性欲差异的合理出口}（高欲一方、或伴侣暂时不便时的正常调节），也是\textbf{了解自身身体的正当方式}。只有当它变成\textbf{强迫性、明显影响生活或关系}时才需要关注。双方可坦诚沟通，让它成为关系的一部分而非秘密。

\section{红旗信号}
\begin{itemize}[leftmargin=2em,itemsep=0.3em]
  \item ED 可能是\textbf{心血管病的首诊信号}——尤其新发 ED 伴胸痛/气促/运动耐量下降，\textbf{尽快就医}。
  \item 持续 3 个月以上、或伴晨勃明显减少、或突然发生——找\textbf{泌尿外科/男科}。
  \item 伴性欲明显下降/情绪低落——可能叠加\textbf{抑郁/激素}问题，一并评估。
\end{itemize}

\section{本章小结}
\begin{itemize}[leftmargin=2em,itemsep=0.3em]
  \item \textbf{叫 ED，不叫"阳痿"}——去污名化，持续 3 个月以上才需评估。
  \item \textbf{ED 常是健康信号}（心血管预警），不是"面子问题"。
  \item \textbf{有一线药}（PDE5 抑制剂），需处方；硝酸酯类同服者禁用。
  \item \textbf{伴侣的支持要点}：不指责、一起就医、当健康问题、亲密不等于插入。
  \item \textbf{自慰是正常的性表达}，不等于不忠；只有强迫性才需关注。
  \item \textbf{红旗}：新发 ED 伴胸痛/气促、持续 3 个月以上、伴性欲明显下降——就医。
\end{itemize}

\chapter{男性生殖系统}
\label{chap:anatomy}

理解男性的性功能与生殖健康，需要先把解剖结构与生理功能弄清楚。本章分两部分：先说明"身体里各是什么、在哪里"，再说明"它们如何工作"，最后归纳几个常见疑问与需要就医的信号。这些是理解后文勃起障碍、早泄、前列腺、不育等问题的基础。

\section{解剖结构}

男性生殖系统分为外生殖器与内生殖器。外生殖器包括阴茎和阴囊；内生殖器包括生殖腺（睾丸）、输送管道（附睾、输精管、射精管及部分尿道）和附属腺体（精囊、前列腺、尿道球腺）。

\subsection*{外生殖器}

阴茎由两根阴茎海绵体和一根尿道海绵体构成：阴茎海绵体左右并列、位于背侧，是勃起的主要结构；尿道海绵体位于腹侧，尿道贯穿其全长，前端膨大为阴茎头，后端膨大为尿道球。海绵体由许多小梁和腔隙组成，与血管相通，腔隙充血时阴茎即变粗、变硬而勃起。阴茎皮肤薄而柔软、富有伸缩性，向前游离延伸形成包绕阴茎头的包皮，包皮前端围成包皮口；阴茎头尿道外口与包皮相连的皮肤皱襞即包皮系带。

阴囊是包裹睾丸的皮肤囊，分为左右两半，各容纳一侧睾丸，通过舒缩调节囊内温度——生精所需温度略低于体温，体外位置正提供这一"降温"。

\subsection*{内生殖器}

\textbf{睾丸}位于阴囊内，左右各一，为稍扁的卵圆形器官，是产生精子和分泌雄激素的生殖腺。睾丸表面覆以浆膜和白膜，后缘白膜增厚形成睾丸纵隔，其结缔组织发出许多小隔，将睾丸分成许多小叶，每小叶含 1--4 条高度盘曲的精曲小管。精曲小管是精子发生场所，管壁由生精上皮构成，上皮下方有基膜及肌样细胞，后者收缩有助于精子排出；精曲小管间为睾丸间质，其中的间质细胞（Leydig 细胞）分泌睾酮。

\textbf{附睾}紧贴在睾丸后上方，是精子成熟和暂时贮存的场所。附睾头部接受来自睾丸的输出小管，尾部连接输精管。精子在附睾内停留约 15--25 天，逐步获得向前运动和受精能力；附睾液也参与精子的营养与成熟。

\textbf{输精管}起自附睾尾部，沿睾丸后缘上行成为精索的一部分，再经腹股沟、盆部向前下方行抵前列腺。进入前列腺前，精囊的输出管加入输精管，二者汇合成射精管穿入前列腺并开口于前列腺部尿道。输精管近端黏膜有纤毛，末端膨大的壶腹部则兼作精子贮存。

\textbf{前列腺}位于膀胱颈下方，包绕尿道起始部，重约 20 g，环绕尿道和两条射精管，由左右叶及前、后、中叶组成。它分泌弱碱性液体，含酸性磷酸酶、纤维蛋白溶解酶、钙等，参与精液的凝固与液化。

\textbf{精囊}位于前列腺后上方，左右各一，分泌量大（构成精液的主要部分），其液体呈碱性并含果糖等，为精子提供能量并中和阴道酸性环境。

\textbf{尿道球腺}很小，位于尿道膜部后下方，射精前分泌黏液润滑尿道；尿道全程还分布有少量尿道小腺体。

\section{生理功能}

\subsection*{精子发生与成熟}

精子发生（spermatogenesis）在睾丸精曲小管内进行，受神经内分泌的严密调控：精原细胞经有丝分裂增殖、分化，再经减数分裂和精子形成（变态），最终生成成熟精子，整个过程约需 64--75 天。

具体而言，精原细胞在青春期后于垂体促性腺激素作用下开始增殖：一部分保持原始状态（A 型），另一部分（B 型）发育为初级精母细胞；初级精母细胞经染色体复制、同源染色体联会与交换后，完成第一次成熟分裂形成两个次级精母细胞；次级精母细胞再迅速完成第二次成熟分裂，形成核型为 23X 或 23Y 的精子细胞。

在精曲小管中生成的精子尚未具备受精能力，须进入附睾后停留约 15--25 天逐步成熟，其形态、膜通透性、代谢和运动方式都随之改变——运动方式由原地摆动，经转圈状运动，最终变为螺旋式的向前运动，并获得受精能力。精子进入女性生殖道后，还需经"获能"才具备穿透卵子的能力；众多精子中最终只有一个与卵子结合，融合形成受精卵后，卵子外围随即形成屏障阻止其他精子进入。

\subsection*{激素调节：下丘脑-垂体-睾丸轴}

睾丸分泌的雄激素主要是睾酮，成年男性每日分泌约 4--9 mg，50 岁后逐渐减少。睾酮的作用包括：促进男性性器官发育和副性征出现、维持性欲与性行为、促进肌肉蛋白质合成、促进骨骼生长与钙磷沉积、促进红细胞生成。

这一过程由下丘脑-垂体-睾丸轴调节并存在负反馈：下丘脑分泌促性腺激素释放激素（GnRH），经垂体门脉到达腺垂体，促使其分泌卵泡刺激素（FSH）和黄体生成素（LH）；FSH 作用于精曲小管、LH 作用于间质细胞，共同驱动生精与睾酮合成。当血中睾酮浓度升高时，反过来抑制下丘脑和垂体，使 GnRH、FSH、LH 分泌减少，从而维持睾酮水平相对稳定。

\subsection*{勃起与射精}

勃起是复杂的神经反射过程，基本反射中枢在脊髓骶段，同时受意识调控。性刺激（视觉、听觉、触觉或想象）经大脑和脊髓传至阴茎，使阴茎海绵体小动脉扩张、血窦充血膨胀；白膜限制其过度扩张，静脉回流受阻，阴茎因而增大变硬。年轻健康男性勃起时，动脉血流量可明显增高。

射精也是有序的多阶段反射：先由附睾、输精管平滑肌顺序收缩，将精子排入后尿道；继而坐骨海绵体肌和球海绵体肌收缩，将精液经尿道射出体外。

\subsection*{精液的组成}

精液由精子和精浆两部分组成。精浆主要由精囊和前列腺的分泌液构成（约占 90\% 以上），另含少量附睾和尿道球腺分泌液；其作用包括运送精子、中和阴道酸性、提供能源物质，以及参与精液的凝固与液化。

\section{几个常见疑问}

\subsection*{关于阴茎大小}

疲软状态下个体差异很大，但\textbf{功能不受尺寸影响}——能否正常勃起并完成性交才是关键。女性阴道的大部分（中后段）对触觉刺激并不敏感，快感主要来自前 1/3 及阴蒂，因此"尺寸"对满意度的影响远小于沟通、前戏与技巧。多数"尺寸焦虑"来自色情内容的筛选效应和同伴比较，与实际体验无关。

\subsection*{关于晨勃与随机勃起}

晨勃、无性念头时的随机勃起，都是\textbf{血管-神经系统的正常自检信号}，说明勃起机制完好，是健康的标志而非异常。青春期起随机勃起的频率会变化，属正常生理，不必刻意压制（焦虑反而使其更易发生）。临床上也可用多种仪器客观记录夜间勃起，用于鉴别心理性与器质性勃起障碍。

\subsection*{关于梦遗}

精液积累到一定程度会自然排出，梦遗（夜间射精）是正常排出方式之一。青春期较常见，随年龄增长和性生活规律而减少；\textbf{频率高低均属正常，不影响健康，也不"伤肾"}。只有当伴随明显不适或突然显著改变时，才需留意（如排除前列腺炎等）。

\subsection*{关于睾丸大小}

成年男性睾丸大小存在一定范围（大致纵径 4--5 cm、横径 2.5--3 cm），在此范围内一般视为正常。睾丸过小者可能提示精曲小管发育不良，影响生精与生育力，可伴勃起或射精力量欠佳，应及时就医评估。

\section{性交中的意外与就医信号}

\begin{itemize}[leftmargin=2em,itemsep=0.3em]
  \item \textbf{阴茎折断（白膜撕裂）}：剧烈性交或体位急转时偶发，典型表现为"啪"的一声加剧痛、迅速肿胀瘀青；若尿道口出血、血尿或排尿困难，提示尿道也受损，应尽快泌尿外科急诊（多可手术修复，延误易留瘢痕与弯曲）。
  \item \textbf{包皮嵌顿}：包皮过长者上翻后卡住无法复位，逐渐肿胀疼痛，需及时复位，拖延可致缺血。
  \item \textbf{睾丸挤压伤}：会阴部受压可致剧痛；若剧痛、持续肿大、伴恶心呕吐，需就医排除\textbf{睾丸扭转}（急症，一般需在 6 小时内处理）。
  \item \textbf{尿道或皮肤损伤}：过剧摩擦或用器具不当可致尿道口或阴茎皮肤损伤、水肿。
  \item 极度疲劳、饱食或饮酒后立即剧烈性交会增加心血管负担。
  \item \textbf{任何性交后持续的疼痛、肿胀、出血或排尿异常，都应及时就医}，不必"忍一忍就过去"。
\end{itemize}

\section{本章小结}

男性生殖系统由睾丸（生精与分泌睾酮）、附睾（精子成熟与贮存）、输精管与射精管（输送）、前列腺/精囊/尿道球腺（分泌）及阴茎/阴囊（外生殖器）构成，其功能受下丘脑-垂体-睾丸轴调控。阴茎大小、晨勃、梦遗等大多是正常生理，不必焦虑；性交中偶发的疼痛、肿胀、出血则应及时就医。理解这些基础，后文谈勃起障碍、早泄、前列腺、不育时才有据可依。


\chapter{前列腺健康：三个常被混为一谈的问题}
\label{chap:prostate}

前列腺是男性\textbf{泌尿与生殖系统的核心腺体}——核桃大小，位于膀胱下方、紧贴直肠前方，环绕尿道。它的功能是为精液分泌约 30\% 的碱性液体（含柠檬酸、酸性磷酸酶、纤维蛋白溶酶等），帮助精液液化、中和阴道酸性。正因为它\textbf{位置特殊（尿道穿过它、直肠紧贴它）}，它的三类问题——\textbf{前列腺炎、前列腺增生、前列腺癌}——常常\textbf{症状高度重叠}，却被很多人混成一句"前列腺有病"，甚至追问"前列腺炎会不会变成癌"。本章把三件事一次讲清：它们\textbf{病因不同、高发年龄不同、处理路径不同，彼此并不必然转化}。

\section{先厘清：三个概念不要混为一谈}

\begin{tabularx}{\textwidth}{l X X X}
\toprule
& \textbf{前列腺炎} & \textbf{前列腺增生 BPH} & \textbf{前列腺癌} \\
\midrule
本质 & 前列腺/盆底\textbf{炎症、疼痛、排尿刺激} & 前列腺\textbf{体积增大}压迫尿道 & 前列腺腺体\textbf{细胞异常增生} \\
高发年龄 & \textbf{15--50 岁}（30 岁前后最多） & \textbf{50 岁后}渐增，70 岁约八成有病理改变 & \textbf{50 岁后}高发，65 岁后最多 \\
典型表现 & 会阴/下腹/睾丸\textbf{疼痛}、尿频尿急、射精痛 & \textbf{排尿费力、尿线细、夜尿、尿不尽} & 早期\textbf{多无症状}；晚期可排尿异常、骨转移痛 \\
是否"性病" & \textbf{多数不是}（仅少数细菌型可经性传播） & 否 & 否 \\
会互相转化吗 & \textbf{没有证据}显示三者必然互转 & 增生\textbf{不等于}癌 & 癌\textbf{不等于}已转移，多数生长慢 \\
\bottomrule
\end{tabularx}

\vspace{0.5em}
\noindent\textbf{关键认知}：排尿困难既可能来自\textbf{增生（压迫）}，也可能来自\textbf{炎症（刺激）}，还可能来自\textbf{癌症}——所以"尿频尿急"不能直接判断是哪一种，\textbf{必须结合年龄、疼痛性质、体检与必要检查}。

\section{前列腺炎：最常见、最易被“抗生素一治了之”}

\textbf{它是什么}：前列腺或周围盆底组织出现\textbf{炎症样疼痛 + 排尿/性功能症状}。按国际 NIH 分类，绝大多数（约九成以上）是\textbf{III 型慢性盆腔疼痛综合征（CPPS）}——\textbf{找不到明确细菌感染}，更像盆底肌、神经、免疫的综合问题；只有少数 I/II 型是\textbf{真正的细菌性}（尿培养能查到菌，多见于年轻人急性发作）。

\textbf{典型表现}：会阴、下腹、腰骶、睾丸、射精管的\textbf{疼痛或不适}（久坐加重）；尿频、尿急、尿后滴沥；射精痛或血精；反复"像尿路感染却查不出菌"。

\begin{itemize}[leftmargin=2em,itemsep=0.3em]
  \item \textbf{别一上来就"抗生素包治"}：非细菌型（占多数）用抗生素\textbf{常无效}，盲目长疗程反而带来肠道菌群紊乱。先做\textbf{尿常规+培养、前列腺液/精液检查、必要时性病筛查}——\textbf{有菌才用抗生素，且有明确疗程}。
  \item \textbf{久坐是明确的加重因素}：每坐 40--60 分钟起身活动；\textbf{避免长时间骑行/压迫会阴}。
  \item \textbf{生活方式（对多数人有实质帮助）}：规律运动、减重、减少酒精/咖啡因/辛辣刺激（会加重盆底充血）、温热坐浴、\textbf{盆底肌放松训练}（注意：是"放松"而非一味收紧）。
  \item \textbf{它不是"性病"、不必然导致不育}：仅少数细菌型可能影响精液；\textbf{别因为"前列腺炎"就默认自己不育或得了性病}。
  \item \textbf{心理因素会放大疼痛}：焦虑→盆底肌更紧→更痛，形成循环；\textbf{正规男科/泌尿外科评估 + 必要时心理支持}，比"换十家医院"更有效。
  \item \textbf{就医路径}：先\textbf{泌尿外科/男科}；反复不愈可找\textbf{盆底康复专科}。
\end{itemize}

\section{前列腺增生（BPH）：50 岁后的"排尿问题"，不等于癌}

\textbf{它是什么}：前列腺随年龄\textbf{体积增大}（50 岁约一半、70 岁约八成有病理增生），增大的腺体\textbf{压迫尿道}→排尿困难。它是\textbf{年龄相关的良性过程}，\textbf{和前列腺癌没有因果关系}——增生不等于癌，癌也不是由增生"变"来的。

\textbf{LUTS（下尿路症状）自查——这些信号要留意}：
\begin{itemize}[leftmargin=2em,itemsep=0.3em]
  \item \textbf{储尿期}：尿频、尿急、\textbf{夜尿增多}（每晚 2 次及以上要重视）。
  \item \textbf{排尿期}：排尿费力、尿线变细、排尿中断、尿末滴沥、尿不尽感。
  \item \textbf{急性尿潴留}（突然排不出尿、小腹胀痛）——\textbf{是急症}，需及时导尿。
\end{itemize}

\textbf{处理路径（由轻到重，需遵医嘱）}：
\begin{itemize}[leftmargin=2em,itemsep=0.3em]
  \item \textbf{轻度}：观察 + 生活方式（睡前限水、减少咖啡因/酒精、别"忍尿"）。
  \item \textbf{药物（一线）}：$\alpha$ 受体阻滞剂（如坦索罗辛，松弛尿道平滑肌、起效快；副作用体位性低血压、逆行射精）；5$\alpha$-还原酶抑制剂（非那雄胺/度他雄胺，缩小体积，起效需 3--6 个月，副作用性欲下降/勃起障碍/射精量减少——\textbf{会使 PSA 值下降约一半}，判读 PSA 时须告知医生在用）。
  \item \textbf{手术/微创（中重度或反复潴留、血尿、结石）}：TURP（经尿道前列腺电切，长期"金标准"）、激光剜除（HoLEP/绿激光）等；选择取决于前列腺大小、技术条件与合并症。
  \item \textbf{术后常见}：逆行射精（精液反流入膀胱，\textbf{常见且多可接受}，不影响勃起与快感）、偶尔尿失禁（多随时间改善）。
\end{itemize}

\noindent\textbf{重要}：BPH 患者\textbf{同样要按年龄做前列腺癌筛查}——增生的"大"和癌的"恶性"是两码事，\textbf{别因为"已经增生/做过手术"就以为排除了癌}。

\section{前列腺癌：早期多无症状，关键是"分级 + 分期"而非只看 PSA}

\textbf{它是什么}：前列腺腺体细胞异常增生、失控生长。它是男性\textbf{最常见的恶性肿瘤之一}，但绝大多数生长\textbf{缓慢}，很多患者\textbf{终生不会因它死亡}——这正是"过度诊断/过度治疗"争议的核心。

\textbf{筛查与诊断（需个体化评估）}：
\begin{itemize}[leftmargin=2em,itemsep=0.3em]
  \item \textbf{PSA（前列腺特异抗原）}：血液指标，\textbf{升高不等于癌}——炎症、增生、射精后、骑车、直肠刺激都能抬高；\textbf{正常也不等于排除癌}。\textbf{PSA 是"提示"而非"确诊"}。
  \item \textbf{数字直肠指检（DRE）}：触诊前列腺有无硬结，\textbf{但很多癌 DRE 摸不到}。
  \item \textbf{多参数 MRI（mpMRI）}：PSA 升高或可疑时，MRI 引导的穿刺比"盲穿"更精准。
  \item \textbf{穿刺活检 + Gleason 评分 / Grade Group}：\textbf{真正的分级依据}——Grade Group 1（低危）到 5（高危）决定后续策略。
  \item \textbf{筛查的争议}：美国 USPSTF 对 55--69 岁建议\textbf{个体化决策}（权衡筛查获益 vs 过度诊断、尿失禁/性功能障碍风险）；有家族史、PSA 早期升高者更倾向筛查；\textbf{一般不建议 70 岁以上新开始筛查}。
\end{itemize}

\textbf{治疗策略（按风险分层，核心是"别一刀切"）}：
\begin{itemize}[leftmargin=2em,itemsep=0.3em]
  \item \textbf{低危（Grade Group 1--2、PSA 低）}：\textbf{主动监测（active surveillance）}——定期 PSA + MRI + 必要时复查穿刺，\textbf{暂不手术/放疗}，避免过度治疗带来的尿失禁·性功能障碍。这是当前低危的主流策略。
  \item \textbf{中高危}：根治性前列腺切除术（开放/腹腔镜/机器人）或根治性放疗（外照射 $\pm$ 内分泌）$\pm$ 短期内内分泌治疗；选择取决于年龄、合并症、肿瘤分期分级、对排尿/性功能的保留需求。
  \item \textbf{转移性}：以\textbf{内分泌治疗（雄激素剥夺 ADT）}为主，联合新型内分泌药（如阿比特龙、恩扎卢胺）或化疗（多西他赛）；近年\textbf{PSMA 靶向（镭-223 等）}与免疫治疗为后线带来新选择。
  \item \textbf{副作用管理}：ADT 常见\textbf{性欲下降、骨质疏松、血脂/血糖变化、潮热、疲劳}——需长期监测骨密度、血脂、体重，必要时补充钙·维生素 D、规律运动、抗骨质疏松治疗。
\end{itemize}

\section{高频误区辨析}

\begin{tabularx}{\textwidth}{l X}
\toprule
\textbf{误区} & \textbf{事实} \\
\midrule
"前列腺炎 = 性病" & \textbf{多数不是}；仅少数细菌型可由性传播（淋病/衣原体）。 \\
"前列腺炎会变癌" & \textbf{没有证据}支持这种必然转化。 \\
"增生 = 癌" & \textbf{增生是良性}体积增大，癌是细胞恶性，两码事。 \\
"PSA 高就是癌" & \textbf{高不等于癌}，炎症/增生/射精都能抬高；正常也不排除癌。 \\
"前列腺癌 = 绝症" & \textbf{大多数生长慢}，早期发现多可长期控制，甚至可主动监测。 \\
"前列腺越大越危险" & 大小与"是否恶性"\textbf{无直接关系}，关键是分级/分期。 \\
"年轻不会得前列腺癌" & 少见但可能；\textbf{有家族史者}风险更高。 \\
\bottomrule
\end{tabularx}

\section{本章小结}
\begin{itemize}[leftmargin=2em,itemsep=0.2em]
  \item \textbf{三件事不是一件事}：炎症（年轻、疼痛）/ 增生（50 岁后、排尿难）/ 癌（50 岁后、多无症状）。
  \item \textbf{前列腺炎}：多数非细菌型，\textbf{别盲目抗生素}；久坐是加重因素，运动·减重·放松更实在。
  \item \textbf{增生 BPH}：排尿症状自查；药物（坦索罗辛/非那雄胺）或手术（TURP/激光）；\textbf{增生不等于癌}，仍要筛查。
  \item \textbf{前列腺癌}：PSA 是提示非确诊；\textbf{关键看 Grade Group + 分期}；低危首选\textbf{主动监测}，别一刀切。
  \item \textbf{就医}：泌尿外科/男科；反复不愈找盆底康复；癌按风险分层个体化。
  \item \textbf{红旗（立刻就医）}：突然排不出尿（急性尿潴留）、反复血尿/血精、进行性排尿困难、骨痛（疑转移）。
\end{itemize}

\vspace{1em}
\noindent\textbf{要点}：前列腺的三件事——炎症（疼）、增生（胀）、癌（恶性）——\textbf{病因、年龄、处理都不同，彼此不必然转化}；\textbf{尿频尿急先分清是哪一种，再谈治疗}，别自己"对号入座"。

\chapter{男性 STI 自查：男性特有的线索与窗口期}
\label{chap:male-sti}

性传播感染（STI）的总论、传播途径与检测逻辑见 A 册“性传播感染详解”。本章聚焦\textbf{男性一侧最容易被忽略的自查线索、检测窗口期，以及"看着像性病、其实不是"的常见误诊}——因为很多 STI 早期\textbf{没有症状}，而有些"看着吓人"的表现（如冠状沟一圈小疙瘩）其实是\textbf{正常变异}。自查的目的是\textbf{早发现、早检测}，\textbf{不是自己给自己下诊断}。

\section{认知前提：症状只是线索，确诊靠检测}

\begin{itemize}[leftmargin=2em,itemsep=0.3em]
  \item \textbf{很多 STI 潜伏/无症状}：梅毒早期、衣原体感染、HPV 感染、HIV 急性期都可"看着没事"——\textbf{无症状不等于没得}。
  \item \textbf{一个症状可能对应多种病}：同一个"尿道口流点东西"，淋病、衣原体、念珠菌性龟头炎都可能——\textbf{必须靠培养/核酸/血清学区分}。
  \item \textbf{有些"症状"根本不是性病}：珍珠状丘疹、系带旁丘疹、皮脂腺异位、毛囊炎——\textbf{看着像湿疣，其实是正常/良性}。
  \item \textbf{自查能做的}：观察\textbf{部位、形态、时间线、是否伴随其他症状}，这些是就医时最有用、也最容易被"不好意思说清"的关键信息。
\end{itemize}

\section{男性特有的自查线索（按部位）}

\begin{tabularx}{\textwidth}{l X X}
\toprule
\textbf{部位} & \textbf{可能线索} & \textbf{需留意的病种（举例）} \\
\midrule
尿道口 & 流脓（黄/绿）、滴白、排尿灼热/疼痛 & 淋病（流脓明显）、衣原体（滴白、轻症） \\
龟头/包皮 & 红斑、小水疱、无痛性溃疡、菜花样赘生物、瘙痒脱屑 & 疱疹（水疱→溃疡）、梅毒（无痛硬下疳）、HPV（疣）、念珠菌（红、痒、白屑） \\
腹股沟 & 淋巴结肿大、压痛 & 梅毒、疱疹、软下疳、腹股沟肉芽肿（少见） \\
阴囊/睾丸 & 单侧肿痛（附睾炎）、硬结 & 淋球菌/衣原体上行、腮腺炎性睾丸炎、睾丸扭转（急症） \\
全身 & 发热、关节痛、躯干/手掌皮疹 & 梅毒二期、HIV 急性期、反应性关节炎（反应性） \\
\bottomrule
\end{tabularx}

\noindent\textbf{关键区分——两种"溃疡"}：
\begin{itemize}[leftmargin=2em,itemsep=0.3em]
  \item \textbf{硬下疳（梅毒一期）}：单个（或少数）\textbf{无痛、质硬、边界清}的溃疡，通常\textbf{不痒不痛}，数周后自行消退（但病菌仍在，进入二期）。
  \item \textbf{生殖器疱疹（HSV）}：\textbf{成簇小水疱}→破溃成\textbf{疼痛性浅溃疡}，常\textbf{反复发作}，发作前有"发麻/灼热"前驱感。
  \item 二者都\textbf{可查、可控}，但\textbf{病因不同、处理不同}——这正是"必须检测"的原因。
\end{itemize}

\section{检测窗口期：从暴露到可检出}

\textbf{窗口期} = 感染后到检测方法能测出标志物的时间。太早测可能"假阴性"——没到窗口期，不是没得。以下为\textbf{大致范围}，具体以医生判断为准：

\begin{tabularx}{\textwidth}{l X X X}
\toprule
\textbf{病种} & \textbf{典型潜伏期（出现症状} & \textbf{能检测的大致窗口} & \textbf{常用检测} \\
\midrule
淋病 & 约 2--7 天 & 有症状即可查；核酸较早期可检出 & 尿道拭子/尿 NAAT \\
衣原体 & 约 1--3 周 & 症状后或 1--2 周 & 尿/拭子 NAAT \\
梅毒（一期） & 硬下疳约 2--6 周 & 硬下疳出现后约 1--2 周血清学转阳 & RPR/TRUST + TPPA \\
疱疹 HSV & 约 4--7 天（2--20 天） & 溃疡期拭子核酸；血清滞后 & 拭子 PCR（发作时最佳） \\
HPV（疣） & 约 2 周--8 个月 & 疣体出现后可查（多凭形态） & 视诊 + 必要时醋酸白/病理 \\
HIV & 急性期约 2--4 周 & 四代 Ag/Ab 约 2--4 周；4 周阴性基本排除 & 四代 Ag/Ab（或核酸检测更早） \\
乙肝 & 约 1--6 个月 & HBsAg 约 1--2 个月 & HBsAg + 抗体系列 \\
\bottomrule
\end{tabularx}

\vspace{0.5em}
\noindent\textbf{实用结论}：高风险行为后\textbf{不必"马上测"，也不必"等太久"}——按上表"能测的窗口"安排一次检测，\textbf{不确定就间隔 4--6 周复查一次}；\textbf{出现症状则随时可查}（不必等窗口期）。

\section{易误诊：看似性病实则不是}

\begin{tabularx}{\textwidth}{l X}
\toprule
\textbf{表现} & \textbf{更可能是}（而非性病） \\
\midrule
冠状沟一圈排列整齐的"小珍珠" & \textbf{珍珠状阴茎丘疹}（正常变异，非性病、不传染、无需治疗） \\
系带旁一排小丘疹 & \textbf{系带旁丘疹}（正常变异） \\
龟头/包皮白色小颗粒、无痛 & \textbf{皮脂腺异位}（正常变异） \\
阴囊/大腿根部红点、痒、与摩擦相关 & \textbf{毛囊炎/股癣/湿疹}（真菌或刺激，非 STI） \\
龟头红、痒、白色分泌物、异味 & \textbf{念珠菌性包皮龟头炎}（真菌，需处理但非性传播） \\
\bottomrule
\end{tabularx}

\noindent\textbf{区分要点}：性病（如湿疣）多\textbf{形态不规则、逐渐增多、边界不清}；正常变异多\textbf{形态规整、长期不变、对称}。\textbf{不确定就查}，别对着手机图片"对号入座"。

\section{男性应做哪些检测}

\begin{itemize}[leftmargin=2em,itemsep=0.3em]
  \item \textbf{基础}：中段尿/首段尿 + 尿道拭子\textbf{核酸检测（NAAT）}（淋球菌、衣原体）——\textbf{最常用、最灵敏}的筛查。
  \item \textbf{血清学}：梅毒（RPR/TRUST + TPPA）、HIV（四代 Ag/Ab）、乙肝（HBsAg 等）。
  \item \textbf{有溃疡/疣}：病灶拭子\textbf{PCR/核酸}（梅毒、HSV）或视诊 + 必要时病理（HPV）。
  \item \textbf{有全身症状}（发热/皮疹/关节痛）：医生会结合\textbf{血常规、相关血清学}综合判断。
  \item \textbf{提醒}：\textbf{检测前 24 小时别排尿}（中段尿标本更准）、\textbf{别自行用药}（会干扰结果）、\textbf{别做剧烈前列腺按摩/直肠指检前}留尿（可能污染）。
\end{itemize}

\section{高频误区}

\begin{tabularx}{\textwidth}{l X}
\toprule
\textbf{误区} & \textbf{事实} \\
\midrule
"没症状 = 没得" & \textbf{很多 STI 无症状}，梅毒/HPV/HIV 早期尤甚。 \\
"一个症状定一个病" & \textbf{同一症状可对应多种病}，靠检测区分。 \\
"冠状沟小疙瘩 = 湿疣" & \textbf{常是珍珠状丘疹}（正常变异），别自己吓自己。 \\
"测过一次阴性就放心" & \textbf{未到窗口期的阴性不可靠}，按窗口期复查。 \\
"淋病能好，就不查其他" & \textbf{可合并衣原体/其他 STI}，建议一并筛。 \\
"吃两天药就不去查了" & \textbf{不规范治疗}易复发/耐药/传播给伴侣。 \\
\bottomrule
\end{tabularx}

\section{本章小结}
\begin{itemize}[leftmargin=2em,itemsep=0.2em]
  \item \textbf{症状是线索，确诊靠检测}——别自己"看图诊断"。
  \item \textbf{男性特有线索}：尿道流脓/滴白、龟头水疱或无痛溃疡、腹股沟淋巴结、阴囊肿痛、发热伴皮疹。
  \item \textbf{硬下疳（无痛、硬） vs 疱疹（成簇水疱、痛、复发）}——都需检测。
  \item \textbf{窗口期}：太早测可能假阴，按上表安排；\textbf{有症状随时可查}。
  \item \textbf{正常变异}（珍珠疹、系带旁丘疹、皮脂腺异位）\textbf{不是性病}，别焦虑。
  \item \textbf{检测}：NAAT（淋/衣）+ 梅毒 + HIV + 乙肝血清学 + 病灶 PCR。
  \item \textbf{就医}：皮肤性病科/泌尿外科/男科；\textbf{伴侣同查同治}，治疗完成前避免无保护性行为。
  \item \textbf{红旗（尽快就医）}：无痛性溃疡、反复发作的水疱、持续淋巴结肿大、发热伴全身皮疹。
\end{itemize}

\vspace{1em}
\noindent\textbf{要点}：男性 STI 自查的核心是\textbf{观察部位·形态·时间线}，再按\textbf{窗口期}去做\textbf{核酸检测 + 血清学}——\textbf{别自己"看图定病"，也别"没症状就放心"}。

\chapter{男性绝育：输精管结扎（vasectomy）}
\label{chap:vasectomy}

输精管结扎是"关掉精子输送通道"的小手术，\textbf{不关激素、不关勃起、不关射精量}。它是所有避孕方式里最省事、最便宜、最接近"一次搞定"的选项之一，也是最容易被误解的一个——很多人以为它是"阉割"，其实完全不是。

\subsection*{先澄清三个最常见的误解}

\begin{center}
\renewcommand{\arraystretch}{1.3}
\begin{tabularx}{\textwidth}{|X|X|X|}
\hline
\textbf{常见说法} & \textbf{实际情况} & \textbf{为什么} \\
\hline
结扎 = 阉割 & 不是。不切睾丸、不切性腺 & 结扎的只是"输精管"，睾丸还在正常工作 \\
结扎后立刻就没精子了 & 不是。需 6 个月或约 14–20 次射精后残留精子排空 & 远端管道和精囊里还存着结扎前的精子 \\
结扎会导致阳痿 & 极少因此发生。多数是心理因素 & 勃起靠血液、靠神经、靠激素，都和输精管无关 \\
结扎后还能轻松恢复 & 可能，但不保证。复通率约 95\%，受孕率约 75\% & 复通比结扎难得多，且可能产生抗精子抗体 \\
结扎会明显减少射精量 & 不会明显改变。精液主体来自精囊、前列腺 & 精子只占精液体积的一小部分 \\
\hline
\end{tabularx}
\end{center}

\subsection*{手术原理：结扎改变了什么、没改变什么}

睾丸每时每刻都在产生精子，这些精子本来经附睾、输精管排出体外。结扎就是\textbf{在输精管上打断这段通道}（结扎、切断或堵塞），让精子排不出去，于是精液里没有精子，女方因此不会怀孕。

关键在于，\textbf{结扎只断了"运输路"，没动"工厂"}：

\begin{itemize}
  \item 睾丸仍持续生精，只是比结扎前\textbf{速度略慢}；
  \item 排不出去的精子会在附睾内被机体\textbf{重吸收}，不会堆积爆炸；
  \item 精液的\textbf{体积、颜色、射精的感觉基本不变}（因为精子只占精液的一小部分，液体主要来自精囊和前列腺）。
\end{itemize}

\subsection*{常见手术方式}

\textbf{输精管结扎术}是我国使用最早、最广泛的男性绝育术。按操作方式可分为：

\begin{center}
\renewcommand{\arraystretch}{1.3}
\begin{tabularx}{\textwidth}{|X|X|X|}
\hline
\textbf{方式} & \textbf{做法} & \textbf{特点} \\
\hline
传统切开结扎 & 阴囊两侧切口，游离并切除约 1 厘米输精管，两残端结扎 & 经典做法，需要时再做吻合复通 \\
\hline
直视钳穿法 & 用皮外固定钳把输精管连阴囊皮肤一起套住固定，一次分开、提出输精管，压挫结扎并用\textbf{筋膜包埋}处理残端 & 全程直视下完成，不用锐器进阴囊，出血少、时间短 \\
\hline
银夹绝育 & 分离出一段输精管，用 V/U 形银夹夹住钳紧，造成管腔闭塞 & 以机械压迫阻截精子 \\
\hline
穿刺电凝 & 特制绝缘针插入输精管腔，高频电流烫伤一小段黏膜，结痂后瘢痕阻塞管道 & 微创，依靠瘢痕闭管 \\
\hline
注射粘堵 & 阴囊皮肤穿刺，将液体粘堵剂注入输精管，凝固后阻塞管道 & "打一针"绝育，怕开刀者接受度高 \\
\hline
可复性栓堵 & 穿刺注入栓堵剂形成有弹性的栓子堵塞管道；需要复孕时取出栓子即可 & \textbf{可逆性最好}，国内外广泛应用 \\
\hline
\end{tabularx}
\end{center}

其中\textbf{可复性栓堵术}和\textbf{粘堵术}可逆性较好，可逆性栓堵在我国处于世界领先水平、已广泛应用。

\subsection*{术后过渡期与继续避孕}

\textbf{这是最容易被忽略、也最关键的一点。}结扎后，输精道远端和精囊里还残留着活精子，所以\textbf{术后不能立刻认为安全}：

\begin{itemize}
  \item 术后一般应\textbf{继续避孕 6 个月或排精 14–20 次}；
  \item 条件许可时做\textbf{精液检查}，确认确实无精子，才能停用其他避孕；
  \item 在此之前坚持用安全套等其他方式；
  \item 为加速残余精子排空，我国早在 1958 年（吴阶平提出）就有术中向输精管远睾端\textbf{低压缓慢灌注杀精药液}的做法，已列入节育手术常规（如 0.01\% 醋酸苯汞、新洁尔灭、维生素 C 等，具体由医生决定）。
\end{itemize}

\subsection*{并发症与副作用}

输精管结扎是\textbf{安全、简便的门诊小手术}，较严重并发症罕见。术后短期内约\textbf{半数}男子会出现阴囊局部瘀血、肿胀、疼痛，一般不严重，多在 1–2 周内自行消退。少数人可能出现下列情况：

\begin{center}
\renewcommand{\arraystretch}{1.3}
\begin{tabularx}{\textwidth}{|X|X|X|}
\hline
\textbf{情况} & \textbf{发生时间} & \textbf{处理要点} \\
\hline
出血 / 阴囊血肿 & 多在术后 24 小时内 & 促凝、抗感染、理疗促吸收 \\
\hline
感染（切口、精索、生殖道） & 术后早期 & 急性期广谱抗生素 + 热敷理疗；化脓者切开引流 \\
\hline
痛性结节 & 术后 3 个月以后，术部压痛 & 消炎、局部封闭、理疗；长期不愈可手术切除结节 \\
\hline
附睾郁积症 & 术后 3 个月以后，阴囊胀痛、附睾肿大压痛 & 阴囊托 + 理疗；长期不愈者吻合复通 \\
\hline
性功能障碍（性欲减退、阳痿、早泄等） & 时间不定 & \textbf{多为心理因素}；先排除器质性疾病，以心理疏导为主 \\
\hline
\end{tabularx}
\end{center}

\textbf{需要立即就诊的信号}：发热（38℃ 以上）、伤口渗血或有脓性分泌物、严重疼痛或肿胀。

\subsection*{恢复生育：输精管复通术}

需要恢复生育力时，可施行\textbf{输精管吻合术}（俗称再通术）。因为结扎或粘堵只是打断/损伤了极少一段管道，剩余部分再接通的可能性较大；栓堵术则把栓子取出即可复通。若睾丸、附睾功能正常、远端管道通畅，多可恢复生育。

\begin{itemize}
  \item 目前普遍采用\textbf{显微外科技术}，减少组织损伤、提高修复成功率；
  \item 复通率通常可达\textbf{约 95\%}，但长期随访\textbf{实际受孕率约 75\%}；
  \item 精子重新出现\textbf{早则 1 个月、晚则半年到 1 年}，早期无精子不代表手术失败；
  \item 结扎术后\textbf{50\% 以上}可能出现血清\textbf{抗精子抗体}，它是干扰受孕的因素之一；
  \item 若复通后仍不怀孕，应做\textbf{系统的生育力检查}，不要只怪结扎。
\end{itemize}

\textbf{提醒}：复通比结扎难得多、且结果不保证——所以\textbf{在还没有准备好、或未来 5–10 年很可能还要孩子时，请慎重选择结扎}。

\subsection*{术前准备}

\begin{itemize}
  \item 术前由专业人员向受术者及家属\textbf{讲清原理与预期}，对预防术后心理障碍极有帮助；
  \item 术前 1 日淋浴、更换清洁内衣；
  \item 术前刮除阴毛，用肥皂水、清水洗净阴囊部。
\end{itemize}

\subsection*{术后注意}

\begin{itemize}
  \item 术后当天静息，48 小时内避免体力劳动和剧烈运动；
  \item 使用阴囊托 48 小时（若戴着更舒服可延长）；
  \item 保持伤口清洁干燥 48 小时；避免撕扯、揉搓、刺激伤口；术后 24 小时不要洗澡，非洗不可务必保持伤口干燥；
  \item 一周后可恢复性生活，无不适应可提早，感觉不佳则不勉强；
  \item 术后一段时间内仍可能使伴侣怀孕，请继续用其他方式避孕，直到精液检查确认无精子（有条件时在约 20 次射精后送精液检查）；
  \item 疼痛或不适可服用常规止痛剂；
  \item 阴囊区疼痛、肿胀、局部青紫属\textbf{正常现象}，不必过度担忧（出现前述"立即就诊信号"除外）。
\end{itemize}

\subsection*{本章小结}

\begin{itemize}
  \item 结扎 = 断精子运输路，\textbf{不是阉割}，不影响激素、勃起和射精量。
  \item \textbf{术后仍会怀孕一段时间}：需避孕 6 个月或排精 14–20 次，并以精液检查确认无精子。
  \item 是安全的门诊小手术；瘀血、肿胀、疼痛多可在 1–2 周自行消退。
  \item 少数可出现血肿、感染、痛性结节、附睾郁积、心因性性功能障碍，均\textbf{可处理}。
  \item 复通率约 95\%、受孕率约 75\%，且可能伴抗精子抗体——\textbf{复通不保证}，做之前想清楚。
  \item 近期内很可能还要孩子？\textbf{别急着结扎}。
\end{itemize}


% =================================================
% C 册补充专题：男性作为性侵受害者（呼应 A 册《性侵犯与强奸》）
% 现代循证撰写
% =================================================

\chapter{男性作为性侵受害者（专科补充）}

\section{先放下“男人就该扛住”的观念}
性侵不只是女性会遇到的事。\textbf{男性、男性之间的、以及异性恋男性被女性施加的}性侵都存在，
只是长期被"男人要主动、要硬、要能行"的刻板印象遮住，导致男性受害者\textbf{更少求助、更晚被发现}。
对男性来说，最大的障碍往往不是疼痛，而是\textbf{羞耻与被质疑}（"一个大男人怎么会""你当时怎么不反抗"）。
请把这两句先收下：\textbf{你是受害者，与你的性别、体力、是否"想发生"都无关；求助不是软弱。}

\section{男性特有的医学要点}
\begin{itemize}
  \item \textbf{生殖道 / 泌尿道损伤}：尿道或阴茎撕裂、血尿、射精痛，可能伤及尿道海绵体；出现"排尿困难、血尿、阴囊肿痛"要尽快到\textbf{急诊 / 泌尿外科}。
  \item \textbf{STI 暴露}：与 A 册一致——HIV（72 小时内 PEP 最有效）、梅毒、淋病、HPV、肝 B 等，需检测与随访。
  \item \textbf{意外致他人妊娠}：若为阴道 / 肛交且无保护，\textbf{伴侣}可能怀孕——这是男性受害者常被忽略、却又必须面对的一环。
  \item \textbf{会阴疼痛、盆底肌紧张}：部分男性事后出现会阴痛、勃起或射精不适，多与盆底肌痉挛 + 心理因素相关，可治疗。
\end{itemize}

\section{就医时如何表述}
\begin{itemize}
  \item 直接说：\textbf{"我可能遭遇了性侵，想检查一下伤情和感染风险，请帮我做相关检查。"}——这样说即可，不必展开细节。
  \item 可要求：体格检查 + 分泌物/尿液取样 + 血检（HIV/梅毒/乙肝）+ 必要时 STI 随访 + 伴侣妊娠评估的指引。
  \item 若不愿走法律程序，可只体检取证、保留证据，之后再决定是否报案；身体安全永远第一。
  \item 找\textbf{男性医生}或\textbf{同性别医护人员}陪同检查，能明显降低不适感。
\end{itemize}

\section{心理与社会支持}
\begin{itemize}
  \item 男性受害者同样可能出现 PTSD、闪回、对亲密关系退缩、性欲骤变（可能升高也可能降低，都正常）。
  \item 支持渠道：心理科 / 创伤聚焦咨询（TF-CBT、EMDR）、男性权益或心理援助热线；必要时伴侣共同咨询。
  \item 一句给男性的话：\textbf{"扛住"不是勇敢，"说出口"才是；你不必为别人的暴力负责。}
\end{itemize}

\section{本章小结}
男性也是性侵受害者，路径与 A 册一致：安全→就医（泌尿外科/急诊）→STI/损伤/伴侣妊娠评估→心理支持。
男性专科要点：生殖道与泌尿道损伤、会阴/盆底疼痛、致伴侣妊娠风险。
开不了口时，一句话"我可能遭了性侵，想检查"就够；找同性别医护陪同、按自己节奏来。




\clearpage
\chapter{男性不育：病因、检查与处理}
\label{chap:infertility}

男性不育指\textbf{在未采取避孕措施、规律性生活 1 年以上仍未使配偶受孕}。它是\textbf{夫妻双方的问题，不全是"男人的事"}——约 1/3 纯男方因素、1/3 纯女方因素、1/3 双方或不明。本章系统讲男方一侧的常见病因与处理思路，具体检查和治疗以\textbf{男科/生殖医学科}为准。

核心观念先立住：
\begin{itemize}[leftmargin=2em,itemsep=0.3em]
  \item \textbf{多数男性不育可查、可治或可通过辅助生殖解决}——不是"判了死刑"。
  \item \textbf{第一步是精液常规分析}（禁欲 2--7 天后，至少 2 次）+ 男方体检——这是所有检查的起点。
  \item 病因分\textbf{生成障碍}（精液质量差）/ \textbf{输送障碍}（精液排不出去）/ \textbf{功能、免疫、全身}三大类，处理路径完全不同。
\end{itemize}

\section{生殖系统肿瘤}

能影响生育的男性生殖系统肿瘤包括阴茎癌、尿道癌、睾丸肿瘤、精囊囊肿、附睾肿瘤，以及前列腺增生和癌。要点如下：

\begin{itemize}[leftmargin=2em,itemsep=0.3em]
  \item \textbf{阴茎癌}：多起于阴茎头和冠状沟，与包茎、个人卫生差、不洁性交（尤其梅毒）相关。早期因勃起时疼痛或出血便已中止性交，就诊时多已偏晚。佩罗尼氏病（阴茎海绵体纤维斑块硬结）会使勃起时向一侧弯曲、严重者不能性交、常合并阳痿。
  \item \textbf{睾丸肿瘤}：临床常见病，94\% 来源于生殖细胞，以精原细胞瘤与畸胎癌常见。单侧肿瘤本身通常不影响生育，但\textbf{腹膜后淋巴结清除术}可能引起射精障碍；化疗可致睾丸萎缩、脱发，需较长时间才能逐渐恢复生精。
  \item \textbf{尿道肿瘤（尤其癌）}：少见。发病不久即充血疼痛、出血、排尿痛、性交痛并伴尿道狭窄，主要因性交困难或不能性交而影响生育。
  \item \textbf{精囊、精索、附睾肿瘤}：精囊以囊肿为主、多为单侧，生育并非完全不可能；精索肿瘤几乎全为胚胎性、多良性，若精道通畅影响不大；附睾肿瘤虽少见，但会影响精子质量、数量和输送。
\end{itemize}

\section{生殖系统炎症}

男性生殖器官炎症与泌尿系统炎症关系密切、常互为因果。

\begin{itemize}[leftmargin=2em,itemsep=0.3em]
  \item \textbf{前列腺炎（尤其慢性）}：成年男性常见，病程某阶段可出现阳痿、早泄、神经衰弱而影响生育；出现这些症状时\textbf{首先要排除前列腺炎症}。
  \item \textbf{精囊炎 / 前列腺炎（细菌或病毒）}：都可抑制生精功能。即使无症状，精液中细菌也能引起精细胞分解、粒细胞中毒和精子凝集。
  \item \textbf{睾丸炎}：常见原因是附睾炎、腮腺炎、梅毒。流行性腮腺炎所致者多发生在腮腺炎一周后、多为单侧；\textbf{双侧受累可因睾丸萎缩导致不育}，可见生精上皮严重损害。淋球菌感染可致睾丸功能丧失、附睾部分或完全阻塞。
  \item \textbf{附睾炎（尤其慢性）}：常与慢性前列腺炎、精囊炎同时或先后发生，可因结核杆菌、淋病双球菌或非特异性感染引起；炎性病变使附睾管腔瘢痕变窄，可能发展到完全闭塞，双侧则致不育。
  \item \textbf{包皮炎 / 龟头炎}：常见病，病期因性交困难而不育，一般病程较短、愈后影响不大，但\textbf{须警惕淋病与梅毒}在阴茎头的表现。
  \item \textbf{无症状生殖系感染}：如生殖道支原体属感染，可造成不成熟精子早期脱落。
\end{itemize}

\section{精子输送障碍}

指精道本身正常、但精液不能射入配偶生殖道的情况，男女双方都可能有原因。男方一侧的常见原因（可对症治疗）：

\begin{itemize}[leftmargin=2em,itemsep=0.3em]
  \item \textbf{先天性畸形}：常见为一侧或双侧输精管缺如/闭锁，也有半附睾缺失。
  \item \textbf{炎症}：生殖系统急、慢性炎症均可致管腔狭窄或阻塞（附睾、输精管、精阜、前列腺等）；慢性炎中\textbf{特别注意结核杆菌}的特异性炎症，也须警惕性病所致炎症。
  \item \textbf{肿瘤}：多为良性（附睾腺瘤样瘤、平滑肌瘤、囊肿、精索腺瘤样瘤）及前列腺癌、阴茎癌等，可"内堵外压"影响精子通过。
  \item \textbf{外伤}：阴囊或会阴外伤可损伤附睾、输精管，引起管腔部分或全部断裂，愈合瘢痕也影响通过。
  \item \textbf{其他}：巨大腹股沟疝、睾丸扭转等也可引起精道梗阻。
\end{itemize}

\section{遗传性病因}

下列遗传性疾病可导致不育（具体核型与表现以遗传科评估为准）：

\begin{center}
\renewcommand{\arraystretch}{1.3}
\small
\begin{tabularx}{\textwidth}{|X|X|}
\hline
\textbf{综合征} & \textbf{关键特征} \\
\hline
先天性睾丸发育不全（Klinefelter） & 染色体异常伴睾丸发育不良，典型的遗传性不育 \\
\hline
XYY 综合征 & 男子中约占 0.2\%，身材高大，可无精或少精，性激素基本正常或升高 \\
\hline
Turner 相关（46XY 为主） & 身材矮小、耳位低、盾状胸、肘外翻、心血管畸形；隐睾、造精障碍、雄激素低 \\
\hline
性逆转综合征 & 外表为男但睾丸小、阴毛稀少，核型 46XY \\
\hline
单纯支持细胞综合征 & 约占 3\%；睾丸缺生精细胞、无精子，促性腺激素升高但睾酮正常 \\
\hline
纤毛呆滞（Kartagener）综合征 & 近亲婚配发病率高；内脏反位、支气管扩张、鼻窦炎；精子尾部纤毛运动失常 \\
\hline
Fanconi 综合征 & 出生身材矮小、骨骼畸形（拇指缺损、指趾畸形）、生殖器发育不良；全血细胞减少 \\
\hline
Prader-Willi 综合征 & 婴幼儿期肌张力低下、肥胖、矮小、智力低下、性腺发育不良、隐睾、小阴茎 \\
\hline
多腺性内分泌不全症 & 常染色体隐性遗传；恶性贫血、重症肌无力、肺纤维症、慢性活动性肝炎、性腺功能低下 \\
\hline
Laurence-Moon / 相关综合征 & 或与孕期放射/化学因素致下丘脑障碍有关；肥胖、视网膜色素变性、智能低下、多指（趾）、尿道下裂 \\
\hline
Leriche 综合征（髂总动脉狭窄/闭塞） & 下肢与男性生殖器供血不足；早期勃起困难、下肢疲劳肌萎缩，晚期性功能完全丧失 \\
\hline
\end{tabularx}
\end{center}

\section{免疫因素：抗精子抗体}

精子是一种\textbf{自身抗原}。正常情况下睾丸存在坚固的\textbf{血睾屏障}，保护精子抗原不接触免疫系统，故不引起自身免疫反应。当屏障被打破（精道阻塞、睾丸炎症和损伤、附睾或附属性腺感染，\textbf{尤其是输精管阻塞}），精子抗原与免疫系统接触，可诱发体液免疫和调理免疫反应，产生抗精子抗体。

\begin{itemize}[leftmargin=2em,itemsep=0.3em]
  \item 抗精子抗体可存在于男方、女方或双方；不育夫妇中，\textbf{免疫性继发不育占约 15\%–30\%}。
  \item 人类抗精子抗体主要有两种：\textbf{精子凝集抗体}（致精子凝集）、\textbf{精子制动抗体}（使精子丧失活动力）。
  \item 不育与抗体的确切因果关系仍不完全清楚，需结合其他检查综合判断。
\end{itemize}

\section{其他因素：烟酒·射线·营养·内分泌·呼吸道}

\begin{itemize}[leftmargin=2em,itemsep=0.3em]
  \item \textbf{吸烟}：烟雾中的物质可诱发细胞畸变、影响精子 DNA 与蛋白质合成，使畸形精子比例升高；每天 20 支以上者畸形发生率显著增高。吸烟也是动脉硬化（血管源性阳痿）的危险因素之一。
  \item \textbf{饮酒}：酒精是性腺毒素，长期过量可致血清睾酮降低、性欲减退、精子畸形、阳痿；每天烈酒 250 mL 以上持续 2–5 年可致严重阳痿。有上述表现者应下决心戒酒。
  \item \textbf{放射线}：对生精细胞敏感。小剂量或偶照后生精功能多可恢复；\textbf{大剂量或长期无防护照射可致不可逆损伤或染色体畸变}。应严格防护并定期体检。
  \item \textbf{维生素与微量元素}：缺乏维生素 A、B、C、E 均可影响精子发生；锌缺乏可致不育，补锌对某些不育有效；铁、镁、锌等也影响精子存活率、活动度和穿透宫颈黏液能力。治疗中适当补充 A、B、E、C 是必要的。
  \item \textbf{甲状腺功能低下}：男性可出现性功能低下、性成熟推迟、副性征落后、性欲减退、阳痿和睾丸萎缩，影响生育。
  \item \textbf{慢性呼吸道疾病}：慢性支气管炎/支气管扩张患者（Kartagener、男性鞭毛异常症）可无精或少精，但睾丸体积与 FSH/LH 多正常——因精子尾部纤毛运动异常，属"与呼吸道疾病相关的无精/少精"。
\end{itemize}

\subsection*{治疗思路}

治疗主要针对\textbf{提高精液质量}、\textbf{消除影响因素}（感染、免疫、精索静脉曲张等）、以及\textbf{解除精道梗阻}：

\begin{itemize}[leftmargin=2em,itemsep=0.3em]
  \item \textbf{药物治疗}：内分泌治疗适用于性腺功能低下、少精子症、精索静脉曲张术后及特发性不育；也处理生殖道感染、免疫性精子活动障碍、高催乳素血症、精液不液化、性功能障碍及其他全身疾病。
  \item \textbf{手术治疗}：先决条件是\textbf{睾丸仍具生精功能}；范围包括纠正外生殖器发育反常、勃起功能障碍、改善生精功能及解除精道梗阻。
  \item \textbf{辅助生殖技术}：自然受孕困难时的重要选项（详见相关章节）。
\end{itemize}

\section{中医视角（参考）}

中医典籍所述"五不男"，见于明·万全《广嗣纪要·择配篇》，指男子生殖器官先天发育畸形或其他病变影响生育，即\textbf{天、漏、键、怯、变}五种：

\begin{itemize}[leftmargin=2em,itemsep=0.3em]
  \item \textbf{天}：天宦的简称，男子生殖器短小若无，不能生育；
  \item \textbf{漏}：遗精、滑精，精关不固、常自遗泻而影响生育；
  \item \textbf{键}：阴茎被割除而不能生育；
  \item \textbf{怯}：阳事不举（阳痿）影响生育；
  \item \textbf{变}：生殖器官两性畸形影响生育。
\end{itemize}

《辨证录》另载"男子不能生育有六病"：\textbf{一精寒、二气衰、三痰多、四相火盛、五精稀少、六气郁}。

\begin{itemize}[leftmargin=2em,itemsep=0.3em]
  \item \textbf{精寒}：下焦虚寒、命门火衰，精液温度低（"冷如冰铁"），难以使女方受孕；
  \item \textbf{气衰}：脏腑机能不强、肾气不足，肾精生成动力不足；
  \item \textbf{痰多}：与脾、肺有关，痰湿蕴郁致真气不足、精气亏耗；
  \item \textbf{相火盛}：肾阴亏损、虚火亢盛（命门火旺），出现性欲太过、遗精、早泻；
  \item \textbf{精稀少}：即"虚劳精少"，泄精少甚至一二滴，因先天不足或房事不节、劳心过度耗损精气；
  \item \textbf{气郁}：情志郁结、肝气不舒，可致血瘀，造成阳痿、不射精。
\end{itemize}

\section{常用中成药（参考，需辨证）}

\begin{center}
\renewcommand{\arraystretch}{1.3}
\small
\begin{tabularx}{\textwidth}{|X|X|X|}
\hline
\textbf{药名} & \textbf{功效} & \textbf{适用与用法（示例）} \\
\hline
龟龄集 & 补肾壮阳 & 肾气虚弱、阳事不举、婚久不育、腰膝疲软；散剂每次 1.5 g，日 1–2 次 \\
\hline
大造丸 & 补肾滋阴、泻热潜阳 & 肾亏阴虚阳亢、骨蒸劳热、消瘦、不育；蜜丸 9 g，日 2 次 \\
\hline
通关丸 & 滋阴降火、引火归元 & 肾阴亏耗、相火妄动、滑精白浊、头晕耳鸣；每次 9 g，日 2 次空腹服 \\
\hline
参茸卫生丸 & 培补肝肾、益气养血 & 体衰神萎、腰膝酸痛、健忘失眠、梦遗滑精、不育；蜜丸 1 丸，日 2 次 \\
\hline
海马保肾丸 & 滋阴益气、补肾助阳 & 肾气虚寒致阳事痿弱、梦遗滑精、婚久不育；蜜丸 6 g，日 2 次 \\
\hline
金匮肾气丸（桂附地黄丸） & 补益肾阳 & 肾阳不足、腰膝冷痛、阳事不举、小便不利或尿频；蜜丸 1 丸，日 2 次 \\
\hline
知柏地黄丸 & 滋阴降火 & 肾阴亏虚、骨蒸劳热、虚烦盗汗、梦遗；蜜丸 1 丸，日 2 次 \\
\hline
鱼鳔丸 & 滋阴补肾、添精益髓 & 气血虚、肾水不足致阳痿早泄、精亏不育、失眠健忘；蜜丸 1 丸，日 2 次 \\
\hline
锁阳固精丸 & 补肾壮阳、温肾固精 & 肾虚致阴囊潮湿、梦遗滑精、腰膝酸痛；蜜丸 1 丸，日 2 次 \\
\hline
五子衍宗丸 & 滋补肾水、添精补髓 & 阴虚气亏、肾水不足致阳痿、遗泄、不育；蜜丸 1 丸，日 2 次 \\
\hline
六味地黄丸 & 滋补肝肾 & 肾阴不足致不育伴消瘦、头晕、遗精、口干；蜜丸 1 丸，日 2 次 \\
\hline
人参鹿茸丸 & 滋肾益气、补血生精 & 肾精亏损致气血两亏、遗精、腰腿酸软、不育；日 1–2 次，温水或黄酒送服 \\
\hline
三肾丸 & 滋阴益气、补肾壮阳 & 肾气虚弱、阴虚火衰致阳痿、腰腿酸痛、婚久不育；蜜丸 1 丸，日 1–2 次 \\
\hline
大补阴丸 & 补肾水、降虚火 & 肾阴不足、相火偏旺致耳鸣、阴虚发热、盗汗；蜜丸 1 丸，日 1–2 次 \\
\hline
左归丸 & 滋补肾水、调养精血 & 肾精亏损致腰腿疲软、遗精、眩晕眼花、不育；小丸 10 g / 大丸 2 丸，日 2 次 \\
\hline
右归丸 & 温补肾阳 & 先天不足、命门火衰致阳痿、畏寒、小便自遗、婚久不育；1 丸，日 2 次 \\
\hline
全席丸 & 温肾壮阳、固精益气 & 肾阳亏损致精神衰惫、头眩耳聋、遗精盗汗、面色萎黄、婚久不育；1 丸，日 2 次 \\
\hline
\end{tabularx}
\end{center}

\textbf{提示}：中成药须\textbf{辨证施治}，虚寒与虚热、肾阴与肾阳用药截然不同，切勿自行对号入座，请在中医师指导下使用。


\section{辅助功法与按摩（参考）}


方法一：

（1）仰卧位：在腹部用摩腹法操作约10分钟，重点在下腹部，并按揉关元穴，性选用气海穴，女方则选子宫穴，治疗应使腹部深处发热为宜。

（2）仰卧位：以手掌按揉大腿内侧，反复3遍，力量稍重。再选穴位按揉，女性用三阴交穴，男性用太溪穴，稍用力按揉至酸胀并保持约2分钟，然后擦大腿内侧，以热透人深部为宜。

（3）俯卧位：沿整个背腰脊柱内侧，用手掌揉推数遍，力量稍重。再分推腰部10遍，然后按揉腰部肾俞穴，以点按出酸胀的感觉后，缓缓施力，保持按揉3～5分钟。

（4）俯卧位：用手掌分别横擦腰部肾区，斜擦骶部八骨穴，均以热透入身体深部为佳。再施用拿风池颈项法和拿肩井法，共约2分钟。

男性不育患者可于早晨起床前和晚上入睡前各行一次指摩会阴急脉法。其方法为侧卧，舌舐上腭，意守丹田，双腿屈曲，用右手中指指端逆时针方向旋转按摩会阴50～100次。继而仰卧，再以双手中指指端同时逆时针方向旋转按摩急脉50～100次。急脉在腹股沟（俗称大腿根）中点，正置阴毛之边际处。

方法二：

男性不育患者还可以练金刚坐、闭气按摩法。

（1）金刚坐法：盘腿而坐，右足跟抵会阴穴，左足跟抵右足解溪穴，涌泉贴靠在右小腿附阴穴，顺项提顶，竖起脊梁，腰不可软弱。身不可倚靠。两手置扶于两腿血海穴，闭目垂帘，去除心中杂念，待气定神敛后，行闭气按摩法。

（2）闭气按摩法：以鼻吸气闭之，用两掌相搓擦极热，急分两手上下搓摩后腰两边软处，右足跟随手下搓与会阴穴发生碰撞，一面徐徐放气从鼻出，周而复始24遍为1度。可行1～3度。然后变金刚坐法。用上述操练法，使肾脏与生殖系统的废物迅速得到排除和贯通，同时强化呼吸，故亦称“贯通补气法”。

（3）收功法：依次觉察从头至足放松3遍，由“功能态”恢复到正常生理状态。继而缓缓摇头晃腰，左右、前后的摆动脊柱、慢慢舒臂伸脚，搓搓手脸，下坐收功。

注意事项：

1.练功初期有流汗流出；脊背局部发胀的感觉。随功夫的深入，有“腰肾如汤煎”，会阴穴有温热气之感。

2.练功要循序渐进，时间要掌握在20分钟左右，室内应自感不冷，搓摩时要注意命门穴。

3.练功期间要节制性生活，可促进练功效果提高并保证精子成熟。

方法三：

推拿对男性不育是一种很好辅助治疗手法。推拿力量要因人而异，用力均匀，争取先轻后重再轻三个步骤。

（1）束腹法：仰卧位。以两手四指对置在患者脊柱正中线两旁，两手将腹两侧抱拢，由后向腹前摩动，经后腰、侧腹至腹正中线止，反复作束腹式弧形摩动3～5分钟。

束腹时应使两手掌心与皮肤表面贴紧；束腹至腹正中线，两手将腹正中的皮肤肌肉捏紧，并向上拿提；操作时有腰肌向前牵拉感，有腰部轻松及增益精神的感觉；有腹胀者慎用。

（2）下腹横摩法：平卧位。以手指或两手四指掌侧并置于患者下腹部左或右侧，即髋骨内上缘前方约平脐四指处，向腹正中平行摩动，反复横摩5～10分钟。

操作时手掌心微向上抬高。手指掌侧平放摩动；横摩时用力宜稍重；便秘患者左侧横摩用力加重。

（3）命门穴按摩：俯卧位。以食指背屈或拇指掌侧揉腰部之命门穴2.5分钟。再以命门穴为中心，一手掌心置其上，作旋转按摩1～2分钟。

指揉时应轻揉有力，揉动时嘱患者腰背肌放松；操作时防止擦伤皮肤。

方法四：

精子畸形症的按摩：提睾功法和阴囊搓擦法，这套功法不但对治疗精子畸形症有效果，对其他精子、精液异常、性功能障碍患者也都有作用。

（1）提睾功法：作用：能提高脏腑的生理功能，起到固肾益气、壮阳补精作用。功法：站立，两足分开与肩同宽，膝盖微曲，站立宁神。两手由下往上抄，掌心朝上，两手随势上托提气。意念由涌泉穴拔气上提，经两足内侧至阴廉，复会阴，将会阴元气上提入腹，至气海、下脘即可。自涌泉至下脘，吸1/2的气（不要吸足）接着慢慢吐气，同时将舌放平，两手掌翻转，中指相对，掌心皆朝下。随着呼气下推。只吐1/4，足趾要抓地，以逆式呼吸法。吸气时收腹，提会阴。用意念将气由涌泉，经足内踝直提至会阴入腹，由下丹田上至中脘，入膻中。两手掌亦随之提到膻中。身体要向前倾30度。练时要通过意念来带动气功的运动。初练时，只可引气至中脘或鸠尾。待功夫练熟了，可提至膻中或天突，甚至可拉到百会穴。初练时，切不可勉强，以防冲气出弊。继之，再反掌。渐渐下推，气由膻中逐节下送至涌泉穴。意念之气由涌泉穴入地3尺，每深1尺，吐一口气。吐气时，足趾放松。

（2）阴囊搓擦法：阴囊搓擦是“提睾功法”的辅助动作，但相当重要。其方法是以手心搓擦阴囊，一手将阴茎上翻，并将两只睾丸推向上面腹股沟内：再以另一手掌心轻轻上下搓揉36次。擦毕，两手交换，以微热为度，每次擦5～10分钟，不要太重，当心擦破囊皮。

方法五：

死精过多的夫妻按摩：

（1）脐旁横摩法。用手掌或食指、中指、无名指指面于脐旁有节律横向抚摩，每分钟120次。

（2）下腹横摩法，用手掌或手指指面于气海、幽门、关元穴有节律横向抚摩，每分钟120次。

（3）揉命门法。用大鱼际、手掌根部或手指吸定命门穴，做轻柔的回旋揉抚，每分钟120～160次。

手足按摩：

（1）掐点手部掌侧命门点；按揉背侧全息穴肾点，下腹点。

（2）点揉足部底侧涌泉穴；背侧公孙穴。

（3）重按手部反应区肾区。生殖区及生殖腺区；掌腕部摩推。

（4）点揉足部反射区生殖区及生殖腺区，肾区；踩足底；擦足心。

\section{本章小结}

\begin{itemize}[leftmargin=2em,itemsep=0.3em]
  \item 不育是\textbf{夫妻双方的问题}，约 1/3 纯男方、1/3 纯女方、1/3 双方或不明——别把责任全压在一方。
  \item \textbf{第一步是精液常规}（禁欲 2–7 天、至少 2 次）+ 男方体检，这是所有检查的起点。
  \item 病因分三大类：\textbf{生成障碍}（精液质量差）、\textbf{输送障碍}（排不出去）、\textbf{功能/免疫/全身}，处理路径不同。
  \item 多数男性不育\textbf{可查、可治或可通过辅助生殖解决}，不是"判了死刑"。
  \item 生活方式影响真实存在：\textbf{烟、酒、放射线、缺乏维生素与微量元素}都会打击生精。
  \item 抗精子抗体可致不育（免疫性继发不育占 15\%–30\%），常见诱因是\textbf{精道阻塞、睾丸/附睾损伤或感染}。
  \item 中成药（龟龄集、五子衍宗丸、金匮肾气丸等）\textbf{须辨证施治}，虚寒虚热用药不同，不要自行对号入座。
  \item 治疗路径：\textbf{药物 → 手术 → 辅助生殖}，选择以男科/生殖医学科评估为准。
\end{itemize}

\end{document}
